\clearpage
\fancyhead[L]{\small CONNOISSANCE DES TEMS.}
\hypertarget{native-542-lower}{}\label{native-542-lower}
\begin{center}\small Göttingische gelehrte Anzeigen. Stück 199. Seite 1978..1987. 1811 December 14.\\[3pt]\rule{.72\linewidth}{.3pt}\end{center}

\emph{Connoissance des tems ou des mouvemens célestes, à l'usage des astronomes et des navigateurs pour l'an} 1812, \emph{publiée par le bureau des longitudes. Juillet} 181c. Paris bei \textsc{Courcier}. 415 Seiten gr. Octav.

Die Tafel für die Längen und Breiten der vornehmsten Oerter der Erde hat bei diesem Jahrgange erhebliche Verbesserungen erhalten; es sind, besonders aus \textsc{Monteiro}'s Ephemeriden von Coimbra und aus den \textsc{Oltmanns}'schen Untersuchungen über die Geographie des neuen Continents über hundert neue Artikel hinzu gekommen, so dass die Anzahl aller jetzt nahe an 1500 beträgt. Auch hat Hr. \textsc{Burckhardt} die Längen von 50 Sternwarten oder sonst durch astronomische Beobachtungen merkwürdi\SourceLeaf{543}ger Punkte in Europa von neuem mit Sorgfalt discutirt; ob indess die Worte il a comparé et calculé de nouveau \emph{toutes} les observations tant anciennes que modernes ganz buchstäblich zu verstehen sind, lassen wir dahin gestellt sein; auf alle Fälle wäre zu wünschen, dass Hr. \textsc{Burckhardt} die Details dieser nützlichen Arbeit umständlicher bekannt machte. \textsc{Bouvard}'s Beobachtungen auf der kaiserl. Sternwarte im Jahre 1809 machen diesmal beinahe die Hälfte der \emph{Additions} aus. Man weiss, wie schätzbar die Sonnen-, Monds-, Planeten- und Sternbeobachtungen sind; wir haben in diesem Jahrgange mit Vergnügen eine Anzahl beobachteter Durchgänge des Polarsterns durch den Meridian bemerkt, nur Schade, dass es immer blos untere Culminationen sind, wobei der vielfache Nutzen, welchen zahlreiche Beobachtungen des Polarsterns, \emph{in beiden Culminationen zugleich}, leisten können, freilich wegfällt. Für Ceres und Pallas, um die Zeit ihrer Opposition, finden sich ziemlich viele Beobachtungen. Eben so zahlreiche Verfinsterungen von Jupiterstrabanten, aber nur drei unvollständig beobachtete Sternbedeckungen, nemlich der Austritt von $\gamma$ Scorpii den 28. Mai $12^{\mathrm u}\,3^{\mathrm m}\,53^{\mathrm s}$ Mittlere Zeit, Austritt von $2\delta$ Tauri den 28. Sept. $9^{\mathrm u}\,43^{\mathrm m}\,0^{\mathrm s}$, Eintritt von $1\delta$ Tauri den 25. Octob. $18^{\mathrm u}\,17^{\mathrm m}\,15^{\mathrm s}2$. --- Leichtes Mittel, Oerter des Mondes näherungsweise zu berechnen, von J. C. \textsc{Burckhardt}. Hr. B. braucht dazu die bekannte Chaldäische Periode von 18 Jahren oder 223 Lunationen, nach deren Verlauf die Argumente der Mondsungleichheiten wieder nahe die vorigen Werthe erhalten, und berechnet die Aenderungen der Länge und Breite, wovon die beträchtlichsten in fünf Tafeln gebracht sind. Aus einer Ephemeride des Mondslaufes für ein gegebenes Jahr kann man so mit sehr weniger Mühe und mit ziemlicher Genauigkeit eine ähnliche Ephemeride für ein um Eine Periode späteres Jahr berechnen. Wir wünschten, dass Hr. \textsc{Burckhardt} zugleich die Aenderung der Horizontalparallaxen auf eine ähnliche Art behandelt hätte. --- Ueber ein neues Mittel die Pendeluhren zu vervollkommnen, von eben demselben. Um den Rost an den Zapfen der Räder zu verhüten, schlägt Hr. \textsc{Burckhardt} vor, sie im Feuer zu vergolden und dann von neuem zu härten. Die wirkliche Ausführung dieses Vorschlags wird am besten lehren können, ob man dabei gewinnt. --- Tafeln für die Aberration, Nutation und Präcession der 36 \textsc{Maskelyne}'schen Fundamentalsterne, von demselben. Mit der Aberration ist zugleich die Solarnutation vereinigt. --- Ueber die Depression des Quecksilbers in den Barometerröhren vermöge der Capillarität, von \textsc{Laplace}. Nach \textsc{Laplace}'s Theorie bewirkt nicht die Capillar-Action des Glases, sondern die einer auch beim sorgfältigen Auskochen des Quecksilbers noch an der Glasfläche zurückbleibenden äusserst feinen Haut von Feuchtigkeit die convexe Oberfläche des Quecksilbers, und damit die Depression desselben in der Röhre. Die Bestimmung der Gestalt der Oberfläche, von welcher die Relation zwischen der Weite der Röhre und der Depression des Quecksilbers abhängt, hat \textsc{Laplace} durch eine mühsame, auf mechanische Quadratur gegründete, Integration der in dem Supplemente zur Mécanique céleste aufgestellten Grundgleichung bestimmt (wobei, nach unserer Meinung, eine etwas weiter getriebene Entwickelung in Reihen vielleicht noch einige Erleichterung verstattet haben würde), und so für den practischen Gebrauch eine hier mitgetheilte Tafel berechnet; wir wünschten, dass die bei dieser Berechnung sich ergebende Höhe des convexen Theils des Quecksilbers auch mit beigefügt wäre. --- Ueber die (in der Monatl. Correspondenz Septemberheft 1808) vom \textsc{Gauss} gegebene Auflösung einer Aufgabe der sphärischen Astronomie, wo aus drei gleichen Höhen bekannter Sterne zugleich die Polhöhe des Orts, der Stand der Uhr und der Fehler des Instruments bestimmt werden, von \textsc{Delambre}. Die Behandlung der Aufgabe, welche Hr. \textsc{Delambre} hier aufstellt, ist im \emph{Wesentlichen} von der von \textsc{Gauss} gegebenen nicht verschieden, hat aber, nach unserer Meinung, an Einfachheit und Concinnität verloren. Hr. \textsc{Delambre} hat das ganze Beispiel einer Beobachtung, womit \textsc{Gauss} damals die Berechnung erläuterte, wieder mit vieler Weitläuftigkeit durchgerechnet; wenn er aber dem von \textsc{Gauss} gefundenen Resultate über den Einfluss der Beobachtungsfehler auf die Genauigkeit der Breitenbestimmung \SourceLeaf{544} \[ \mathrm{d}\varphi = 3{,}8077\Delta - 0{,}2884\Delta' + 3{,}5193\Delta'' \] die Bemerkung beifügt: Cette dernière formule prouve que dans la pratique la méthode n'aurait qu'une exactitude assez bornée: so lässt sich dies unpassende Urtheil nicht anders erklären, als dass er, kaum begreiflicher Weise, übersehen hat, dass $\Delta$, $\Delta'$, $\Delta''$ in Zeitsecunden, und $\mathrm{d}\varphi$ in Bogensecunden ausgedrückt sind. Eben so zeigt der Zusatz, dans une nuit, où l'on pourroit observer trois étoiles à la même hauteur, on verroit très probablement passer au méridien quelque étoile connue qui donneroit la latitude avec moins de peine et plus de précision, et l'heure de la pendule par une simple hauteur, dass Hr. \textsc{Delambre} abermals vergessen hat, dass diese Methode für alle Fälle bestimmt ist, wo man sich auf sein Instrument in Ansehung der absoluten Höhen nicht verlassen kann, und was die \emph{Genauigkeit} betrifft, so ist die Bemerkung ohne allen Grund. --- Es folgt dann noch Einiges über das bekannte, aber in der Ausübung so gut wie ganz unbrauchbare, Problem, aus drei Höhen Eines Sterns zugleich dessen Declinationen, Stundenwinkel und die Polhöhe zu bestimmen. Hr. \textsc{Delambre} kömmt hierauf noch einmal auf die kleine Abhandlung von \textsc{Gauss} zurück, worin die Beobachtung zweier Höhen zweier bekannter Sterne zur Zeit- und Breitenbestimmung vorgeschlagen war, und wovon er in dem vorhergehenden Bande der Connaissance des tems gesprochen hatte. Wir finden in dem, was er darüber und über den Vorzug der Synthese vor der Analyse sagt, einen neuen Beweis unsers bei Anzeige jenes Bandes der Connoissance des tems geäusserten Urtheils, dass Hr. \textsc{Delambre} den Zweck jener Schrift ganz unrichtig aufgefasst habe. Dieser an seinem Orte ganz deutlich ausgesprochene Zweck war, eine \emph{Combination} von Beobachtungen zur Bestimmung der Polhöhe zu empfehlen, die dazu unter manchen Umständen sehr brauchbar ist, und, so viel der Verfasser wusste und bis diese Stunde weiss, dazu in der Allgemeinheit noch nicht vorgeschlagen war, daher er sie eine \emph{neue} Methode nannte. Die \emph{Berechnung} solcher Beobachtungen gründet sich dann, wie \textsc{Gauss} damals gleichfalls zeigte, auf ein Problem, dessen geometrische Auflösung seit \textsc{Tycho}'s Zeiten bekannt ist; wenn daher der Verf., nachdem er diese mit wenig Worten vollständig angedeutet hatte, noch zwei Seiten verwandte, zu zeigen, dass sich dieselbe Auflösung kurz und elegant auch auf rein analytischem Wege finden lasse, so geschah dies nur, weil er glaubte, dass Freunden der Analyse eine solche sich gerade nicht von selbst darbietende und einige Kunst erfordernde Entwickelung angenehm sein könnte, ohne sich träumen zu lassen, dass Jemand dies so auslegen könnte, als ob damit alle geometrische Behandlung, deren Werth bekannt und entschieden genug ist, verdrängt werden sollte. Hrn. \textsc{Delambre}'s Bemerkungen über die gegenseitigen Vorzüge des analytischen und geometrischen Verfahrens vor einander, sind daher, wenn auch meistens völlig gegründet, doch durchaus nicht an ihrem Platze. --- Ueber die verschiedenen, von den Astronomen angewandten, Mittel, die Sonnenfinsternisse zu beobachten, von \textsc{Delambre}. Es wird hier eine merkwürdige Stelle aus \textsc{Appians} Astronomicum Caesareum, gedruckt 1540, angeführt, worin zuerst gefärbte Gläser zu Sonnenbeobachtungen vorgeschlagen werden, obwohl Hr. \textsc{Delambre} aus triftigen Gründen es wahrscheinlich findet, dass \textsc{Appian} seinen Vorschlag selbst auszuführen nicht versucht habe. --- Hierauf folgt ein weitläuftiger, 50 Seiten füllender, Auszug aus Gauss \emph{Theoria motus corporum coelestium}. Hr. \textsc{Delambre} hat den grössten Theil der darin enthaltenen Formeln, meist ohne die Beweise, excerpirt, die numerischen erläuternden Beispiele, obwohl mit geringerer Präcision, als in dem Werke selbst, wieder durchgerechnet, und dies mit hin und wieder eingestreueten Anmerkungen hier abdrucken lassen. So empfehlenswerth eine solche Art, dergleichen Werke zu studiren, ist, so wenig scheinen doch solche Excerpte sich zum Abdruck zu qualificiren. Von den Anmerkungen können wir, wegen des beschränkten Raums, hier nur einige berühren. Die Methode, die Berechnung der geocentrischen Oerter der Planeten durch rechtwinklichte Coordinaten sogleich auf den Aequator zu beziehen, hat Hr. \textsc{Delambre} auf \emph{Ein} Beispiel angewandt, und dabei den Vorzug jenes \SourceLeaf{545}Verfahrens vor dem gewöhnlichen nicht recht einsehen können. Allein das ist ganz gegen den Geist jener Methode, die blos für die Fälle bestimmt ist, wo \emph{viele} geocentrische Oerter berechnet werden sollen. Hätte Hr. \textsc{Delambre}, anstatt Eines Planetenortes, ein Dutzend nach jener Methode berechnet, so würde er dasselbe gefunden haben, was schon so manche andere Rechner fanden, dass man dabei nicht halb so viele Zeit und Mühe nöthig hat, als bei dem gewöhnlichen Verfahren. Die vier Formeln der sphärischen Trigonometrie, wovon in der Theoria motus so vielfacher Gebrauch gemacht ist, hat Hr. \textsc{Delambre} seinerseits auch gefunden, allein ihren Vorzug vor dem gewöhnlichen Verfahren nicht erkannt. Sollen z.\,B. aus zwei Seiten und dem eingeschlossenen Winkel eines sphärischen Dreiecks alle übrigen Stücke bestimmt werden, so hat man nach den neuen Formeln an sechs verschiedenen Stellen der Sinustafeln zusammen 12 Logarithmen aufzusuchen, und dann hat man zugleich eine Controlle der Rechnung und allemal scharfe, nie zweideutige, Resultate; dagegen muss man bei dem gewöhnlichen, von Hrn. \textsc{Delambre} vorgezogenen, Verfahren, wenn man gleichfalls eine Controlle der Rechnung haben will, zusammen 13 Logarithmen an 11 verschiedenen Stellen der Tafel aufsuchen, und erhält dann die dritte Seite durch ihren Sinus, also, wenn sie nahe am rechten Winkel fällt, nicht scharf, ja vielleicht sogar zweideutig, ohne dass sie es im Problem selbst ist. Der Vorzug der neuen Formeln ist daher ganz entschieden, und erheblich genug, wenn man viele dergleichen Operationen zu machen hat, und es ist daher zu verwundern, wie Hr. \textsc{Delambre} ihn hat übersehen können. Einer Berichtigung bedarf der Ausdruck, dessen sich Hr. \textsc{Delambre} S.\,357 in Betreff der zweiten \textsc{Gauss}ischen Auflösung des Problems, aus zwei Abständen eines Planeten von der Sonne, dem eingeschlossenen Winkel und der Zwischenzeit die Elemente zu bestimmen, bedient, dass sie \emph{sehr lange} Entwickelungen erfordere, die der Verfasser nicht gegeben habe, und zum Theil auf ihm eigenthümliche Theorien, die er noch nicht bekannt gemacht habe, gegründet sei. Man sollte hiernach glauben, dass jene Auflösung einer der wichtigsten Aufgaben des ganzen Werks so vorgetragen sei, dass Leser, die nicht vorzüglich geübt sind, gar nicht damit fertig werden, und selbst Kenner doch keine ganz vollständige Einsicht in dieselben erhalten können. Beides ist aber unrichtig. Rec. weiss aus vielfachen Beispielen, dass bei den unbedeutenden, kleinen Entwickelungen, die, wenn man nicht ein durch widerliche Weitläuftigkeit ungeniessbares Buch schreiben will, immer dem Leser überlassen bleiben müssen, selbst Anfänger nirgends Anstoss gefunden haben, und mit den dem Verfasser eigenthümlichen Theorien, deren Entwickelung, hier nicht an ihrem Platze, er sich auf eine andere Gelegenheit vorbehalten musste, hängt in dieser Auflösung nichts zusammen, als die Rechnungsvortheile, die er selbst angewandt hat, um die Hülfstafel zu construiren, und die hierbei durchaus nicht wesentlich sind. Bei der Hauptgleichung in der grossen Aufgabe, die Bahn aus drei geocentrischen Oertern zu bestimmen, hatte \textsc{Gauss} die indirecte Auflösung als vorzüglich bequem empfohlen, aber über die Art, wie dieselbe auszuführen sei, nichts weiter hinzu gefügt, weil theils die dabei anzuwendenden Kunstgriffe an sich bekannt genug sind, theils die in der Theoria gegebene Auflösung des \textsc{Kepler}'schen Problems dabei gewisser Massen als Muster dienen kann. Die Art indess, wie Hr. \textsc{Delambre} die numerische indirecte Auflösung der Hauptgleichung in dem von ihm im grössten Detail wieder durchgerechneten Beispiele angreift, und die, obwohl er selbst sie noch für bequem genug hält, doch mehr als drei Mal zu lang ist, da man mit zwei Versuchen weiter reichen kann, als Hr. \textsc{Delambre} mit sieben, zeigt, dass es doch nicht unzweckmässig sein würde, wenn der Verfasser gelegentlich an einem schicklichen Orte auf die in der indirecten Auflösung anzuwendenden kleinen Kunstgriffe nochmals aufmerksam machte. --- Nach der langen Anzeige der Theoria, bei welcher Hr. \textsc{Delambre} sich nur auf das erste Buch, und den ersten Abschnitt des zweiten, beschränkt hat, folgen noch: Neue Bemerkungen über die Parallaxenrechnung und über die Formeln der Herren \textsc{Olbers} und \textsc{Littrow}, von \textsc{Delambre}, worin derselbe von neuem erklärt, dass er solche Auflösungen von Aufgaben, die sich auf die \SourceLeaf{546}Beziehung der Lage der Punkte im Raume auf drei rechtwinklichte Coordinaten gründen, und wobei die Entwickelungen rein analytisch geschehen, immer für weniger einfach hält, als solche, wobei blos die sphärische Trigonometrie angewandt wird: ein sehr einseitiges Urtheil, mit welchem wenige Mathematiker übereinstimmen werden. Eine Formel, worin die Längenparallaxe durch eine Reihe dargestellt wird, die nach den Sinus der Vielfache des Abstandes des Mondes vom Nonagesimus fortläuft, hatte \textsc{Olbers} dem Hrn. \textsc{Rohde} zugeschrieben; Hr. \textsc{Delambre} reclamirt sie hier als seine Erfindung, die er schon in der Connoissance des tems für 1793 bekannt gemacht habe. Allein eigentlich gehört sie \textsc{Lagrange} zu, der sie schon, noch etwas allgemeiner gefasst, in den Nouveaux Mémoires de l'académie de Berlin a. 1776 p. 231 bekannt gemacht hatte. --- Mittel, um eine Uhr die Sternzeit und mittlere Zeit zeigen zu lassen, von \textsc{Burckhardt}, gründet sich auf das genäherte Verhältniss beider Zeiten, wie 51.79 zu 49.82, welches in einem ganzen Jahre nur 4 Secunden fehlt. Von \emph{demselben}, über den zweiten Cometen von 1737, nach Beobachtungen in China, die in der Monatl. Corresp. bekannt gemacht worden sind. --- Die meteorologischen Beobachtungen auf der Pariser kaiserl. Sternwarte im Jahre 1808, und das Verzeichniss der Mitglieder des Längenbüreau, machen, wie gewöhnlich, den Beschluss dieses Bandes.
